Best Bike Paths is designed and implemented with the purpose of enabling Cyclists to both improve their riding experience and contribute reliable information about bike paths. Path reports and obstacle cues are created by Cyclists (manually or by confirming sensor suggestions) and organized to support practical route searches. The project aims to provide a stable and user-friendly platform through which Cyclists and Visitors can discover better routes, share meaningful field evidence, and keep path quality up to date. This document, also known as the Requirements Analysis and Specification Document (RASD), provides all information needed for a smooth implementation of the BBP system, focusing especially on the interactions between the system and the various actors who will interface with it.

\newpage

\section{Purpose}
\label{sec:purpose}%
The purpose of this document is to present a detailed description of Best Bike Paths.
It is addressed to the developers who have to implement the requirements and could be used as an agreement between the customer and the contractors.\\ 
The document is also intended to provide the customer with a clear and unambiguous description of the system's functionalities and constraints, allowing the customer to validate the requirements and to verify if the system meets the expectations.

\subsection{Goals}
\label{subsec:goals}%
\newcounter{g}
\setcounter{g}{1}
\newcommand{\cg}{\theg\stepcounter{g}}

Below there's a table that lists all the goals of the BBP system:
\begin{center}
    \begin{longtable}{ |l|p{0.9\linewidth}| }
        \hline
        \textbf{ID} & \textbf{Description}                                                                   \\
        \hline
        G\cg        &   Support Personal Cycling Activity Tracking\\
        \hline
        G\cg        &   Improving the Quality and Safety of Cycle Paths\\
        \hline
        G\cg        &  Facilitating the Discovery and Navigation of Suitable Routes\\
        \hline
        G\cg        &   Providing Relevant Contextual Information for Planning\\
        \hline
        G\cg        &  Creating a Trusted Community-Based Path Evaluation System\\
        \hline
        \caption{Goals table.}
        \label{tab:goals_tab}%
    \end{longtable}
\end{center}


\section{Scope}
\label{sec:scope}%
Best Bike Paths (BBP) is a mobile/web platform that helps Cyclists find safer, more effective bike routes and share ground-truth about path quality. Users can record trips (GPS + basic stats), review sensor-suggested obstacles, and publish status reports on path segments. Visitors can search origin–destination routes and see multiple alternatives ranked by a score that combines path status and route effectiveness; when several users report on the same segment, BBP merges their inputs by freshness and confirmations.

\newpage

\subsection{World Phenomena}
\label{subsec:world_phenomena}%
\newcounter{wp}
\setcounter{wp}{1}
\newcommand{\cwp}{\thewp\stepcounter{wp}}
\begin{center}
    \begin{longtable}{ |l|p{0.8\linewidth}| }
        \hline
        \textbf{ID} & \textbf{Description}                               \\
        \hline
        WP\cwp      & An Anonymous User wants to create a new account.   \\
        \hline
        WP\cwp      & A User wants to access their account (Login).   \\
        \hline
        WP\cwp      & A User wants to find a cycling route between point A and point B.   \\
        \hline
        WP\cwp      & A User wants to view the details of a route (e.g., status, weather, elevation gain). \\
        \hline
        WP\cwp      & A Registered User wants to record their activity (distance, speed) while cycling (automatic mode).   \\
        \hline
        WP\cwp      & A Registered User wants to plan a route “at the desk” by drawing it on the map (manual mode).   \\
        \hline
        WP\cwp      & A Registered User wants to save a trip in their personal archive.   \\
        \hline
        WP\cwp      & A Registered User wants to add a route.   \\
        \hline
        WP\cwp      & A Registered User wants to manually report a problem (e.g., pothole, obstacle) on a route.   \\
        \hline
        WP\cwp      & A Registered User wants to confirm/reject a problem report (e.g., pothole) automatically detected by the app.   \\
        \hline
        WP\cwp      & A Registered User wants to view their routes.  \\
        \hline
        WP\cwp      & A Registered User wants to log out of their account (Logout). \\
        \hline
        \caption{World Phenomenas.}
        \label{tab:worldph_tab}%
    \end{longtable}
\end{center}

\newpage
\subsection{Shared phenomena}
\label{subsec:shared_phenomena}%
\newcounter{sp}
\setcounter{sp}{1}
\newcommand{\csp} {\thesp\stepcounter{sp}}
\begin{center}
\begin{footnotesize}
    \begin{longtable}{ |l|p{0.5\linewidth}|l|l|}
        \hline
        \textbf{ID} & \textbf{Description} & \textbf{Controller} & \textbf{Observer} \\
        \hline
        SP\csp  & The User enters textual data (credentials, route names, origin, destination). & User & BBP \\
\hline
    SP\csp  & The User interacts with the interface by pressing buttons (e.g., Login, Save, Publish, Search). & User & BBP \\
    \hline
    SP\csp  & The User starts automatic trip recording. & User & BBP \\
    \hline
    SP\csp  & The User stops automatic trip recording. & User & BBP \\
    \hline
    SP\csp  & The User selects the status of a route in manual mode (e.g., Optimal, Requires maintenance). & User & BBP \\
    \hline
    SP\csp  & The User confirms an automatically detected issue (e.g., pothole or obstacle). & User & BBP \\
    \hline
    SP\csp  & The User rejects or corrects an incorrect automatic detection. & User & BBP \\
    \hline
    SP\csp  & The User selects a specific route from the displayed result list. & User & BBP \\
    \hline
    SP\csp  & BBP displays informational or error messages (e.g., Save completed, GPS error). & BBP & User \\
    \hline
    SP\csp  & BBP displays the interactive map. & BBP & User \\
    \hline
    SP\csp  & BBP renders the recorded route trace on the map. & BBP & User \\
    \hline
    SP\csp  & BBP shows real-time statistics (speed, distance, duration) during recording. & BBP & User \\
    \hline
    SP\csp  & BBP displays the trip summary statistics (average speed, elevation gain, total distance). & BBP & User \\
    \hline
    SP\csp  & BBP shows a ranked list of routes based on the computed score. & BBP & User\\
    \hline
    SP\csp  & BBP displays the aggregated (merged) status of a route. & BBP & User \\
    \hline
    SP\csp  & BBP shows weather data associated with the route (temperature, wind, conditions). & BBP & User \\
    \hline
    SP\csp  & BBP displays a confirmation request for an automatic detection (popup / notification). & BBP & User\\
    \hline
    SP\csp  & The GPS sensor provides coordinates (latitude, longitude) to BBP. & GPS Sensor & BBP \\
    \hline
    SP\csp  & The accelerometer sends acceleration vector measurements (x, y, z) to BBP. & Accelerometer & BBP \\
    \hline
    SP\csp  & The gyroscope provides rotation data (x, y, z) to BBP. & Gyroscope & BBP \\
    \hline
    SP\csp  & The external weather service sends JSON data on temperature, wind, and atmospheric conditions. & Weather API & BBP \\
    \hline

        \caption{Shared Phenomenas.}
        \label{tab:sharedph_tab}%
    \end{longtable}
\end{footnotesize}
\end{center}



\section{Definition, Acronyms, Abbreviations}
\label{sec:definition_acronyms_abbreviations}%
\begin{table}[H]
    \begin{center}
        \begin{tabular}{ |l|l| }
            \hline
            \textbf{Acronyms} & \textbf{Definition}                              \\
            \hline
            RASD             & Requirements Analysis \& Specification Document                      \\
            \hline
            CY             & Cyclist                         \\
            \hline
            ACY              & Anonymous Cyclist                         \\
            \hline
            STG             & Student Group                    \\
            \hline
            BBP             & BestBikePaths                   \\
            \hline
            User            & All CYs and ACYs                           \\
            \hline
            API             & Application Programming Interface                           \\
            \hline
            DAX             & Domain Assumption X                           \\
            \hline
            SPX             & Shared Phenomena X                           \\
            \hline
            WPX             & World Phenomena X                           \\
            \hline
            RX              & Requirement X                           \\
            \hline
         \end{tabular}
        \caption{Acronyms used in the document.}
        \label{tab:acronyms}%
    \end{center}
\end{table}




\section{Revision history}
\label{sec:revision_history}%
\begin{itemize}
    \item \textbf{Version 1.0} - 
\end{itemize}


\section{Reference Documents}
\label{sec:reference_documents}%
\begin{itemize}
    \item Specification Document Assignment
    \item IEEE Standard Documentation For RASD 
\end{itemize}

\newpage

\section{Document Structure}
\label{sec:document_structure}%
The document is organized into six main sections, each focusing on a specific aspect, as summarized below.
\paragraph{Introduction:} The first section outlines the project’s goals and purpose, and provides a brief overview of the global and shared phenomena. It also includes key abbreviations and definitions necessary to understand the problem.
\paragraph{Overall Description:} The second section offers a general overview of the problem, explaining the domain, main scenarios, product and user characteristics, as well as the related assumptions, dependencies, and constraints.
\paragraph{Specific Requirements:} The third section presents a detailed analysis of the system requirements, covering external interfaces, functional aspects, and performance expectations.
\paragraph{Formal Analysis Using Alloy:} The fourth section uses Alloy to perform a formal validation of the model introduced earlier. It reports the results of the executed checks and the verified assertions.
\paragraph{Effort Spent:} The fifth section details the contribution of each group member to the development of this document.
\paragraph{References:} The last section lists all bibliographic sources and materials used during the document’s preparation.

